\documentclass[11pt]{article}
\usepackage[margin=1in]{geometry}
\usepackage[T1]{fontenc}
\usepackage{lmodern}
\usepackage{microtype}
\usepackage{amsmath}
\usepackage{booktabs}
\usepackage{enumitem}
\usepackage{xcolor}
\usepackage{listings}
\usepackage{hyperref}
\hypersetup{colorlinks=true,linkcolor=blue!50!black,urlcolor=blue!50!black}
\lstset{basicstyle=\ttfamily\small,breaklines=true,columns=fullflexible,keepspaces=true}
\newcommand{\code}[1]{\texttt{#1}}
\setlist{nosep}
\sloppy
\setlength{\emergencystretch}{3em}

\title{Bounding the in-memory block DAG in \code{f1r3node-rust}\\
\large Code analysis of issue \#933 and a development plan}
\author{Design note for the F1R3FLY.io node team}
\date{9 November 2026 --- \code{dev} at \code{0ccc315d}}

\begin{document}
\maketitle

\section{Summary}

The claim is correct and the defect is evident in the code. The in-memory
\code{DagState} in \path{block-storage/src/rust/dag/block_metadata_store.rs}
holds seven collections keyed by block hash or height. Six of them ---
\code{dag\_set}, \code{child\_map}, \code{height\_map}, \code{block\_number\_map},
\code{main\_parent\_map} and \code{self\_justification\_map} --- receive an entry on
every insert and have no removal path anywhere in the workspace (outside test
fixtures). Only \code{finalized\_block\_set} is capped. Finalization advances
\code{last\_finalized\_block} and sets per-block flags, but never shrinks the
indices; and at startup \code{BlockMetadataStore::try\_new} rebuilds all six from
every metadata row in LMDB, so memory is proportional to the whole chain
history, never to the live frontier. GitHub issue \#933 (opened 4 Nov 2026)
reports the same structures, with a measured 3{,}493 bytes of live heap per
block and about 37\,MB/h of growth per node on an idle 15-node testnet.

The intuition that finalization should allow the in-memory DAG to be
``reset'' is right in spirit but needs one refinement: the correct shape is a
\emph{window} of heights behind the last finalized block (LFB), not a reset to
the fringe, because several readers legitimately walk below the LFB. The
repository already contains the machinery to make this safe --- the ``restore
horizon'' contract for nodes joined from a Last Finalized State (LFS) snapshot,
under which a block the node does not hold is reported as \code{BlockNotHeld}
and deferred rather than treated as a fault. A windowed DAG makes every node
look like an LFS-restored node whose horizon advances with finality. The plan
below reuses that contract rather than inventing a second one.

\section{What the code does today}

\subsection{The growing structures}

\begin{lstlisting}
// block_metadata_store.rs
pub(crate) struct DagState {
    pub(crate) dag_set: imbl::HashSet<BlockHash>,
    pub(crate) child_map: imbl::HashMap<BlockHash, imbl::HashSet<BlockHash>>,
    pub(crate) height_map: imbl::OrdMap<i64, imbl::HashSet<BlockHash>>,
    pub(crate) block_number_map: imbl::HashMap<BlockHash, i64>,
    pub(crate) main_parent_map: imbl::HashMap<BlockHash, BlockHash>,
    pub(crate) self_justification_map: imbl::HashMap<BlockHash, BlockHash>,
    pub(crate) last_finalized_block: Option<(BlockHash, i64)>,
    pub(crate) finalized_block_set: imbl::HashSet<BlockHash>,
}
\end{lstlisting}

\code{add\_block\_to\_dag\_state} inserts into all of the first six on every
\code{add}. The only remover in the file is
\code{prune\_finalized\_cache\_if\_needed}, which trims
\code{finalized\_block\_set} from 50{,}000 to 25{,}000 entries (and does so in
arbitrary iteration order, not oldest-first; \code{is\_finalized} falls back to
the LMDB metadata row, so this is a cache miss rather than a correctness bug).
A workspace-wide search for \code{.remove(} on the other six fields finds
nothing outside tests.

The design is inherited from the Scala \code{BlockMetadataStore} the file
header cites, which also kept its DAG indices in memory for the life of the
process; the Rust port adds three more per-block maps
(\code{block\_number\_map}, \code{main\_parent\_map},
\code{self\_justification\_map}) as hot-path caches, roughly doubling the
per-block footprint. The measured $\sim$3.5\,KB per block is consistent with
six \code{imbl} HAMT entries holding 32-byte \code{Bytes} keys plus a freshly
allocated \code{imbl::HashSet} per block in \code{child\_map}; the child sets
are the likely dominant term and the profile in Phase~0 should confirm that.

\subsection{What finalization does, and does not, do}

\code{BlockDagKeyValueStorage::record\_directly\_finalized} (called from
\code{finalization\_runner.rs}) walks the held, unfinalized ancestry of the new
LFB, runs the finalization effect, then \code{record\_finalized} marks the
blocks finalized in LMDB, adds them to \code{finalized\_block\_set} and
advances \code{last\_finalized\_block}. Nothing is removed from the six
indices. \code{propagate\_ft\_to\_finalized\_blocks} then scans every
finalized block whose cached fault tolerance is below the new value --- a
second place where cost grows with history rather than with the frontier
(partly mitigated by the \code{ft\_lower\_bound} skip).

\subsection{Startup}

\code{BlockMetadataStore::try\_new} calls \code{store.collect} over the entire
\code{block-metadata} LMDB table and folds every row into a fresh
\code{DagState}. Both the time and the resident memory of a restart therefore
scale with total chain length. Issue \#923 (disk growth) is the on-disk twin;
this note deliberately leaves LMDB contents intact and bounds only the
in-memory view.

\section{Why a reset to the fringe is not quite right}

A \code{KeyValueDagRepresentation} snapshot answers two different kinds of
question: \emph{membership} (``does this node hold block $h$?'') and
\emph{structure} (children, heights, main-parent spine). Several production
readers ask these questions about blocks \emph{below} the LFB:

\begin{description}[leftmargin=1.5em,style=nextline]
\item[Late blocks from lagging validators.] The parent-depth rule
  (\code{validate.rs}, ``ANCHOR: the block's OWN parent frontier, never this
  node's tip'') deliberately allows a block whose parents sit arbitrarily far
  below the validating node's LFB. Validation resolves those parents through
  \code{dag.lookup}, which is gated on \code{dag\_set.contains}. If the parent
  were evicted from \code{dag\_set} while still present in LMDB, \code{lookup}
  would return \code{None}, the dependency gate would park the block in the
  Casper buffer waiting for a block the node already has, and nothing would
  ever release it --- the stall class of \#690 and the \#306 ``not held'' storm,
  self-inflicted. \textbf{Membership must stay truthful for every block the
  node holds, pruned or not.}
\item[Mergeable-channel GC.] \code{mergeable\_channels\_gc.rs} enumerates
  \code{height\_map} from its last swept height up to
  \code{floor.block\_number - (max\_parent\_depth + mergeable\_channels\_gc\_depth\_buffer)}
  and needs \code{children} for each candidate. If heights are evicted before
  the sweep has enumerated them, their mergeable data is never collected. The
  retention depth below the floor must cover this ceiling, and eviction must be
  ordered after enumeration (or the sweep's \code{swept\_height} must be
  persisted so a restart does not lose it; today it is process-local and
  restarts from genesis).
\item[Repeat-deploy scope walks.] \code{walk\_blocks\_in\_scope}
  (\code{proto\_util.rs}, \#776) descends \code{deploy\_lifespan} heights
  below the \emph{proposing block's parents}, reading
  \code{dag.block\_number} from the in-memory index to prune the walk. With
  the own-frontier parent rule those parents can be below the node's LFB.
\item[Silent validators.] \code{latest\_messages()},
  \code{non\_finalized\_blocks()} and the clique oracle follow a validator's
  latest message and self-justification chain; a validator silent for longer
  than the window has a latest message below it.
\item[Block API, history horizon, replay.] \code{topo\_sort},
  \code{get\_max\_height}, \code{find} (truncated-hash search over
  \code{dag\_set}), \code{descendants} (legacy merge path),
  \code{rspace\_history\_horizon.rs}, the replay in \code{initializing.rs}
  and \code{BlockStateManager::is\_empty} (\code{topo\_sort(0, Some(1))}) all
  assume the height index covers the whole chain. \code{genesis\_hash()}
  falls back to \code{height\_map.get(\&0)} when the genesis register is
  unset; the register is written on every bootstrap path in
  \code{initializing.rs}, so the fallback can go.
\end{description}

Fork choice itself (\code{estimator.rs}) starts at the LFB and walks
\emph{children} upward, and the DAG merger scopes blocks between the LFB and
the tips; both live entirely in the retained window and need no change.

\section{Design}

\subsection{Principle}

Keep LMDB complete. Make the in-memory indices a cache over a height window
$[W, \infty)$ where $W$ is the \emph{window floor}, and make every accessor
on \code{KeyValueDagRepresentation} fall through to the LMDB metadata row for a
block not in the window. Everything a walk needs about an old block ---
\code{block\_number}, \code{parents[0]} (main parent), the sender's own
justification, the \code{finalized} flag --- is already in
\code{BlockMetadata}. The one relation that is \emph{not} derivable from a
metadata row is \emph{children}; the design persists a height index to LMDB
and makes children-below-window either a persisted index or an explicitly
unsupported query, depending on the caller.

\subsection{Window floor}

After each successful \code{record\_directly\_finalized},
\begin{align*}
  W &= \mathrm{LFB.block\_number} - D,\\
  D &= \max\bigl(\,\code{max\_parent\_depth} + \code{mergeable\_channels\_gc\_depth\_buffer},\\
    &\qquad\quad \code{deploy\_lifespan},\; \code{dag\_memory\_window\_min}\bigr),
\end{align*}
with \code{dag-memory-window-min} a new \code{casper} configuration key
(suggested default 1{,}000 heights). $W$ is monotone non-decreasing and is
additionally capped at the GC sweep's enumerated height (Phase~3). With the
defaults in \code{defaults.conf} (\code{max-parent-depth = 15},
\code{deploy-lifespan = 50}, buffer 10) the minimum dominates, and at the
observed 5--9 blocks per height a 1{,}000-height window costs roughly
25--30\,MB of index memory --- bounded, and tunable by operators.

\subsection{Accessor semantics after the change}

\begin{center}
\small
\begin{tabular}{@{}p{0.34\textwidth}p{0.26\textwidth}p{0.34\textwidth}@{}}
\toprule
Accessor & In window & Below window ($\mathrm{block\_number} < W$) \\
\midrule
\code{contains} / \code{lookup} & in-memory set, then LMDB row & LMDB row (truthful membership) \\
\code{block\_number}, \code{main\_parent}, \code{self\_justification} & in-memory map & derived from LMDB row \\
\code{is\_finalized} & set, then LMDB & LMDB (unchanged) \\
\code{children} & in-memory map & LMDB \code{children-index}, or \code{None} (see 4.4) \\
\code{topo\_sort}, \code{get\_max\_height} & in-memory \code{height\_map} & LMDB \code{height-index} range scan \\
\code{find} (truncated hash) & in-memory scan & LMDB prefix scan on metadata keys \\
\code{genesis\_hash} & register & register (drop the height-0 fallback) \\
\bottomrule
\end{tabular}
\end{center}

\paragraph{Snapshot semantics.} A representation is a point-in-time snapshot;
falling through to LMDB could reveal a block inserted \emph{after} the
snapshot. Gate the fall-through on the row's \code{block\_number} being below
the snapshot's own $W$ (captured at \code{get\_representation} time). The
residual leak is a late block from a lagging validator that lands below $W$
between a snapshot and its use; such a block is settled-history territory and
the only consumers that could observe it are read-only API calls. Record this
as an accepted deviation in the module docs.

\subsection{Children below the window}

Production callers of \code{children}: the estimator (walks up from the LFB,
always in window), \code{mergeable\_channels\_gc::reaches\_ancestor\_check}
(candidates are at most $D$ below the floor, in window by construction of
$D$), and \code{descendants} (legacy merge path, base is at or above the
LFB). No production caller needs children of a block below $W$. Phase~1
therefore returns \code{None} below the window and documents it; a persisted
children index is listed as an optional follow-up only if the block API's
graph renderers need it.

\subsection{Eviction}

Eviction is by height, under \code{global\_lock.write()}, immediately after
\code{record\_finalized}: for each height $h$ in
\code{height\_map.range(..W)}, remove its hashes from the six maps and remove
the height key. Eviction ignores finalization status --- an unfinalized orphan
below $W$ stays in LMDB and remains reachable through fall-through, and if a
later LFB's ancestry merges it, \code{held\_ancestors} still finds and
finalizes it. The GC sweep keeps its own \code{pending} set, so evicting
orphans from the DAG indices does not lose them for GC. Because \code{imbl}
collections are persistent, eviction is $O(\text{evicted} \cdot \log n)$ and
outstanding snapshots are unaffected.

\section{Development plan}

Each phase is independently mergeable and behaviour-preserving until Phase~3
turns eviction on behind a configuration flag.

\subsection*{Phase 0 --- Reproduce and instrument (1--2 days)}
\begin{enumerate}
\item Land the counting-allocator measurement from \#933 as
  \code{block-storage/tests/dag\_state\_memory\_bound.rs}: 120{,}000
  shard-shaped inserts, assert bytes/block over the last 100{,}000; record the
  3.5\,KB baseline and leave the assertion as a regression bound once
  Phase~3 lands.
\item Add gauges \code{dag.window\_floor} and \code{dag.evicted\_total}
  alongside the existing \code{DAG\_BLOCKS\_SIZE\_METRIC} /
  \code{DAG\_CHILDREN\_INDEX\_SIZE\_METRIC} / \code{DAG\_HEIGHTS\_SIZE\_METRIC}
  in \code{multi\_parent\_casper/snapshot.rs}; these are the soak acceptance
  signals.
\item Confirm with a heap profile which of the six collections dominates.
\end{enumerate}

\subsection*{Phase 1 --- Truthful absence: fall-through accessors (3--5 days)}
\begin{enumerate}
\item Add \code{window\_floor: i64} to \code{DagState} and to
  \code{KeyValueDagRepresentation} (initially \code{i64::MIN}, so nothing
  changes).
\item Rewrite \code{lookup}, \code{contains}, \code{block\_number},
  \code{main\_parent}, \code{self\_justification} and their
  \code{\_metered} / \code{\_unsafe} variants to fall through to
  \code{block\_metadata\_index.get} when the hash is not in \code{dag\_set},
  gated as in 4.3. Reuse \code{block\_metadata\_to\_info} for the derivations.
\item Replace the three production direct-field reads with accessors:
  \code{clique\_oracle.rs:250} (\code{block\_number\_map}),
  \code{mergeable\_channels\_gc.rs:258} (\code{height\_map}, via a new
  \code{heights\_in\_range} accessor), and the size gauges in
  \code{snapshot.rs}. Then narrow the field visibility so the compiler
  enforces it (tests use the existing \code{from\_parts} escape hatch).
\item Drop the \code{height\_map.get(\&0)} fallback in \code{genesis\_hash};
  change \code{BlockStateManager::is\_empty} to test the genesis register or
  \code{dag\_set.is\_empty()}.
\item Make \code{validate\_dag\_state} compute contiguity from $W$ rather than
  from the minimum height, so eviction never trips the warning.
\item All existing tests must pass unchanged; add unit tests for each
  accessor with a block present in LMDB but absent from \code{dag\_set}.
\end{enumerate}

\subsection*{Phase 2 --- Persisted height index and bounded startup (3--4 days)}
\begin{enumerate}
\item Add an LMDB store \code{height-index} ($\code{i64} \to$ set of hashes),
  written in \code{insert\_internal} under \code{global\_lock}, with a
  one-time backfill from the metadata scan when the store is empty at
  startup (same shape as the carrier-index watermark bootstrap).
\item Route \code{topo\_sort} and \code{get\_max\_height} through the
  in-memory map when the range is inside the window and through the LMDB
  index otherwise; \code{find} falls back to a key-prefix scan on
  \code{block-metadata}.
\item Change \code{try\_new} to rebuild only heights $\ge W_0$, where
  $W_0$ is derived from the persisted LFB's block number and $D$; LMDB
  remains the source of truth for everything below. Startup becomes
  $O(\text{window})$.
\end{enumerate}

\subsection*{Phase 3 --- Eviction behind a flag (3--5 days)}
\begin{enumerate}
\item Add \code{casper.dag-memory-window-min} and
  \code{casper.dag-memory-window-enabled} (default off for one release).
\item Implement \code{evict\_below(W)} in \code{BlockMetadataStore} and call
  it from \code{record\_directly\_finalized}'s persist closure after
  \code{record\_finalized}, computing $W$ as in 4.2.
\item Couple to GC: persist \code{GcSweep::swept\_height} in an LMDB meta
  slot and cap $W$ at it, or move enumeration to the persisted height index
  so the sweep no longer depends on in-memory \code{height\_map}. The second
  is cleaner and also fixes the restart-from-genesis re-enumeration.
\item Let \code{finalized\_block\_set} eviction follow $W$ instead of the
  arbitrary 50k/25k trim.
\end{enumerate}

\subsection*{Phase 4 --- Verification (soak, 1 week elapsed)}
\begin{itemize}
\item \textbf{Unit/property.} Property test: for a random DAG and random
  $W$, every accessor returns the same answer with and without eviction
  (fall-through equivalence). Add to \code{casper/tests/fork\_choice/}.
\item \textbf{Lagging validator.} Integration test in
  \code{casper/tests/batch1/}: a validator silent for $> D$ heights proposes
  a block citing parents below $W$; the block must be admitted, not parked
  in the Casper buffer.
\item \textbf{GC.} Mergeable data of blocks evicted from memory is still
  collected; sweep survives a restart.
\item \textbf{LFS joiner and restart.} A node restored from an anchor and a
  node restarted with a populated database both come up with a bounded
  \code{DAG\_BLOCKS\_SIZE} and serve \code{topo\_sort} over pruned heights.
\item \textbf{Soak.} Repeat the \#933 idle testnet: 20\,h, 10 validators,
  5 observers. Acceptance: container anon memory flat after warm-up
  ($< 2$\,MB/h drift), \code{dag\_blocks\_size} bounded by
  $D \times$ blocks-per-height, finalization p95 not worse than \code{dev}.
\item Enable the flag by default once the soak passes; remove the flag one
  release later.
\end{itemize}

\section{Risks and open decisions}

\begin{itemize}
\item \textbf{Snapshot leak (4.3).} Accepted deviation; document it. If a
  reviewer objects, the alternative is to stamp each metadata row with the
  \code{dag\_generation} at insert and gate fall-through on it, at the cost
  of a metadata schema bump.
\item \textbf{Fall-through cost.} Each below-window accessor call is one LMDB
  point read instead of a map lookup. Only lagging-validator validation and
  API calls hit it; the hot paths (fork choice, merge, finalizer) stay in
  window. Confirm with the \code{dag.insert.time} and finalizer timing
  histograms during the soak.
\item \textbf{Adjacent unbounded state, out of scope here.}
  \code{invalid\_blocks\_index} and \code{latest\_messages\_index} are
  materialised with \code{to\_map()} on every \code{get\_representation};
  the equivocation tracker grows with equivocations;
  \code{propagate\_ft\_to\_finalized\_blocks} scans all finalized blocks.
  None grows per block at idle, but each deserves its own ticket.
\item \textbf{Disk (\#923).} This plan keeps LMDB complete by design. Pruning
  LMDB below an archival horizon is a separate decision with
  network-protocol implications (what a node can serve to joiners) and should
  not be bundled with the memory fix.
\end{itemize}

\end{document}
